\documentclass[11pt,reqno]{amsart}
\usepackage{amsmath,amsthm}
\usepackage{mathtools}
\usepackage{latexsym}
\usepackage{amsfonts}
\usepackage{amssymb}
\usepackage[dvipsnames]{xcolor}
\usepackage{graphicx}
\usepackage{float}
\usepackage{graphbox}
\usepackage{textcomp}
\usepackage[colorlinks = true, linkcolor = blue, urlcolor  = blue, citecolor = red]{hyperref}
\usepackage[margin=1in]{geometry}
\usepackage{array}
\usepackage{comment}
\usepackage{appendix}

\newtheorem{theorem}{Theorem}[section]
\newtheorem{proposition}[theorem]{Proposition}
\newtheorem{lemma}[theorem]{Lemma}
\newtheorem{corollary}[theorem]{Corollary}
\newtheorem{definition}[theorem]{Definition}
\newtheorem{question}[theorem]{Question}
\newtheorem{remark}[theorem]{Remark}
\newtheorem{example}[theorem]{Example}

\newcommand{\norm}[1]{\left\lVert#1\right\rVert}
\newcommand{\abs}[1]{\left\lvert#1\right\rvert}
\def\R{\mathbb{R}}
\def\C{\mathbb{C}}
\def\N{\mathbb{N}}
\DeclareMathOperator{\Tr}{Tr}
\DeclareMathOperator{\E}{\mathbb{E}}
\DeclareMathOperator{\Prob}{P}
\DeclareMathOperator{\diag}{diag}

% Draft provenance colors.
\definecolor{RevisionDerived}{RGB}{105,105,105}
\definecolor{RevisionLean}{RGB}{0,105,120}
\definecolor{RevisionNewMath}{RGB}{0,45,150}
\newenvironment{revderived}{\begingroup\color{RevisionDerived}}{\endgroup}
\newenvironment{revlean}{\begingroup\color{RevisionLean}}{\endgroup}
\newenvironment{revnewmath}{\begingroup\color{RevisionNewMath}}{\endgroup}
\newenvironment{revlatest}{\begin{revnewmath}}{\end{revnewmath}}
\newenvironment{revrecent}{\begin{revlean}}{\end{revlean}}
\newenvironment{revearlier}{\begin{revderived}}{\end{revderived}}

\begin{document}
\appendix
\section{Concentration via local Lipschitz estimates}
\label{app:second_approach}

\makeatletter
\@ifundefined{revderived}{\newenvironment{revderived}{}{}}{}
\@ifundefined{revlean}{\newenvironment{revlean}{}{}}{}
\@ifundefined{revnewmath}{\newenvironment{revnewmath}{}{}}{}
\@ifundefined{revlatest}{\newenvironment{revlatest}{}{}}{}
\@ifundefined{revrecent}{\newenvironment{revrecent}{}{}}{}
\@ifundefined{revearlier}{\newenvironment{revearlier}{}{}}{}
\makeatother

We prove concentration for the normalized moments
\[
  f_k(X)=\Tr\left(\left((XX^*)^\Gamma\right)^k\right),
  \qquad
  X\in S_{HS}(d^2,s):=\{X\in M_{d^2,s}(\C):\norm X_2=1\}.
\]
The sphere \(S_{HS}(d^2,s)\) is equipped with its uniform probability
measure, i.e. Euclidean surface measure divided by its surface area.

We use the following norm conventions.  For \(X\in M_{m,n}(\C)\),
\[
  \norm X_2=\bigl(\Tr(X^*X)\bigr)^{1/2}
  =
  \left(\sum_{i=1}^m\sum_{\alpha=1}^n |X_{i\alpha}|^2\right)^{1/2}
\]
is the Hilbert--Schmidt norm, and
\[
  \norm X_\infty
  =
  \sup_{\norm v_{\ell_2^n}=1}\norm{Xv}_{\ell_2^m}
\]
is the operator norm.  For \(A\in M_N(\C)\),
\[
  \norm A_1=\Tr|A|,\qquad |A|=(A^*A)^{1/2},
\]
is the trace norm.  We regard \(M_{d^2}(\C)\) as
\(M_d(\C)\otimes M_d(\C)\), and
\(\Gamma=\operatorname{Id}\otimes T\) denotes partial
transpose on the second tensor factor.

The only external estimates used below are Levy's spherical concentration
inequality \cite[Corollary~5.17]{AubrunSzarek2017} and Aubrun's
off-diagonal partially-transposed Wishart moment estimate
\cite[Lemma~3.4]{Aubrun12}; all other reductions are written out.

\begin{theorem}[Localized concentration of normalized moments]
\label{thm:appendixB-local-concentration}
For every fixed \(k\ge1\) and every \(\lambda>0\), there are constants
\(c,C>0\), depending only on \(k\) and \(\lambda\), such that for all
sufficiently large \(d\), all integers \(s\) satisfying
\[
  \frac{\lambda}{2}d^2\le s\le 2\lambda d^2,
\]
and all \(\varepsilon>0\),
\[
  \Prob\left\{
    \left|
      \Tr\left(\left((XX^*)^\Gamma\right)^k\right)
      -
      \E\Tr\left(\left((XX^*)^\Gamma\right)^k\right)
    \right|
    \ge
    \varepsilon d^{-(2k-2)}
  \right\}
  \le
  C e^{-c d^2}
  +
  C\exp\left(-c\frac{d^4\varepsilon^2}{k^2}\right),
\]
where \(X\) is uniform on \(S_{HS}(d^2,s)\).
\end{theorem}

\subsection{A local form of Levy's lemma}

Let \(\sigma_{n-1}\) denote the uniform probability measure on the Euclidean
unit sphere \(S^{n-1}\subset\R^n\), namely Euclidean surface measure divided
by the surface area of \(S^{n-1}\).  Distances on \(S^{n-1}\) are measured
with the ambient Euclidean norm unless explicitly stated otherwise.

\begin{lemma}[Spherical Levy inequality]
\label{lem:appendixB-spherical-levy}
Assume \(n\ge3\).  Let \(F:S^{n-1}\to\R\) be \(L\)-Lipschitz for the
ambient Euclidean distance, with \(L>0\), and let \(M_F\) be a median of
\(F\) with respect to \(\sigma_{n-1}\).  Then for every \(t>0\),
\[
  \sigma_{n-1}\{x\in S^{n-1}: |F(x)-M_F|> t\}
  \le \exp\left(-\frac{n t^2}{2L^2}\right)
\].
\end{lemma}

\begin{proof}
Aubrun--Szarek's Levy lemma, \cite[Corollary~5.17]{AubrunSzarek2017},
states the following stronger median estimate for the geodesic distance
\(g\) on \(S^{n-1}\): if \(n>2\), if \(F\) is \(L\)-Lipschitz for \(g\),
and if \(M_F\) is a median of \(F\), then, for every \(t>0\),
\[
  \sigma_{n-1}\{x: |F(x)-M_F|>t\}
  \le
  \exp\left(-\frac{n t^2}{2L^2}\right).
\]
The ambient Euclidean distance is dominated by the geodesic distance:
\[
  \norm{x-y}_2\le g(x,y)\qquad(x,y\in S^{n-1}).
\]
Thus an \(L\)-Lipschitz function for the Euclidean distance is also
\(L\)-Lipschitz for \(g\), so the displayed estimate is exactly the stated
estimate.
\end{proof}

The following elementary localization is the form needed below; all
probabilities in the statement are with respect to \(\sigma_{n-1}\).

\begin{lemma}[Localized Levy lemma]
\label{lem:appendixB-local-levy}
Assume \(n\ge3\).  Let \(\Omega\subset S^{n-1}\) be a nonempty measurable
set, and let
\(f:S^{n-1}\to\R\) be measurable.  Assume that \(f\) is
\(L\)-Lipschitz on \(\Omega\), with \(L>0\).  If \(M_f\) is a median of
\(f\), then for every \(t>0\),
\[
  \Prob\{|f-M_f|>t\}
  \le
  2\Prob(\Omega^c)
  +
  2\exp\left(-\frac{n t^2}{8L^2}\right).
\]
\end{lemma}

\begin{proof}
Define the McShane extension of \(f|_\Omega\) by
\[
  F(x)=\inf_{y\in\Omega}\bigl(f(y)+L\norm{x-y}_2\bigr),
  \qquad x\in S^{n-1}.
\]
Then \(F=f\) on \(\Omega\) and \(F\) is \(L\)-Lipschitz on the whole sphere.
Indeed, for any \(x,x'\) and any \(y\in\Omega\),
\[
  f(y)+L\norm{x-y}_2
  \le f(y)+L\norm{x'-y}_2+L\norm{x-x'}_2,
\]
and taking the infimum over \(y\) gives
\(F(x)\le F(x')+L\norm{x-x'}_2\); reversing \(x,x'\) gives the Lipschitz
bound.  If \(x\in\Omega\), the choice \(y=x\) gives \(F(x)\le f(x)\), while
the Lipschitz property of \(f|_\Omega\) gives
\(f(x)\le f(y)+L\norm{x-y}_2\) for every \(y\in\Omega\), hence
\(f(x)\le F(x)\).

Let \(M_F\) be a median of \(F\).  Since \(F=f\) on \(\Omega\),
\[
  \Prob\{|f-M_F|>u\}
  \le
  \Prob(\Omega^c)+
  \exp\left(-\frac{n u^2}{2L^2}\right).
\]
Here the exponential term is exactly Lemma~\ref{lem:appendixB-spherical-levy}
applied to the global \(L\)-Lipschitz function \(F\); the extra
\(\Prob(\Omega^c)\) accounts for the points where \(F\) and \(f\) need not
agree.
The median replacement estimate
\[
  \Prob\{|f-M_f|>t\}
  \le
  2\Prob\{|f-M_F|>t/2\}
\]
is elementary: if \(|M_f-M_F|\le t/2\), the event on the left is contained
in \(\{|f-M_F|>t/2\}\); if \(|M_f-M_F|>t/2\), then
\(\Prob\{|f-M_F|>t/2\}\ge1/2\) by the definition of a median.  Combining the
two estimates gives the claim.
\end{proof}

\subsection{Deterministic Lipschitz estimate}

\begin{lemma}[Trace-power perturbation]
\label{lem:appendixB-trace-power-perturbation}
Let \(N,k\ge1\).  For every \(A,B\in M_N(\C)\),
\[
  |\Tr(A^k)-\Tr(B^k)|
  \le
  k\max(\norm A_\infty,\norm B_\infty)^{k-1}\norm{A-B}_1 .
\]
\end{lemma}

\begin{proof}
The telescoping identity gives
\[
  A^k-B^k
  =
  \sum_{j=0}^{k-1}A^j(A-B)B^{k-1-j}.
\]
For every matrix \(M\), \(|\Tr M|\le\norm M_1\), and the trace norm satisfies
\[
  \norm{RST}_1\le\norm R_\infty\norm S_1\norm T_\infty .
\]
Therefore
\[
\begin{aligned}
|\Tr(A^k)-\Tr(B^k)|
&\le
\sum_{j=0}^{k-1}
\norm{A^j(A-B)B^{k-1-j}}_1       \\
&\le
\sum_{j=0}^{k-1}
\norm A_\infty^j\norm{A-B}_1\norm B_\infty^{k-1-j}                 \\
&\le
k\max(\norm A_\infty,\norm B_\infty)^{k-1}\norm{A-B}_1 .
\end{aligned}
\]
\end{proof}

\begin{lemma}[Partial-transpose moment map Lipschitz bound]
\label{lem:appendixB-moment-lipschitz}
Let \(d,s,k\ge1\).  For \(X,Y\in M_{d^2,s}(\C)\), put
\[
  A_X=(XX^*)^\Gamma,\qquad A_Y=(YY^*)^\Gamma.
\]
Then
\[
\begin{aligned}
&\left|
  \Tr(A_X^k)-\Tr(A_Y^k)
\right|                                                     \\
&\qquad\le
2kd\,
\max(\norm{A_X}_\infty,\norm{A_Y}_\infty)^{k-1}
\max(\norm X_\infty,\norm Y_\infty)\,
\norm{X-Y}_2 .
\end{aligned}
\]
\end{lemma}

\begin{proof}
By Lemma~\ref{lem:appendixB-trace-power-perturbation},
\[
\left|\Tr(A_X^k)-\Tr(A_Y^k)\right|
\le
k\max(\norm{A_X}_\infty,\norm{A_Y}_\infty)^{k-1}
\norm{A_X-A_Y}_1 .
\]
Since \(A_X-A_Y=(XX^*-YY^*)^\Gamma\), partial transpose is an isometry for
the Hilbert--Schmidt norm, and
\(\norm M_1\le d\norm M_2\) for \(M\in M_{d^2}(\C)\), we get
\[
  \norm{A_X-A_Y}_1
  \le d\norm{XX^*-YY^*}_2.
\]
Finally,
\[
  XX^*-YY^*=(X-Y)X^*+Y(X-Y)^*,
\]
and the Hilbert--Schmidt ideal inequality gives
\[
  \norm{XX^*-YY^*}_2
  \le
  \norm{X-Y}_2\norm X_\infty+\norm Y_\infty\norm{X-Y}_2
  \le
  2\max(\norm X_\infty,\norm Y_\infty)\norm{X-Y}_2.
\]
Combining the estimates proves the lemma.
\end{proof}

\subsection{Operator-norm good sets}

\begin{lemma}[Gaussian polar decomposition]
\label{lem:appendixB-gaussian-polar}
Let \(m,n\ge1\), and let \(G\in M_{m,n}(\C)\) have independent standard
complex Gaussian entries.  Put \(R=\norm G_2\) and \(X=G/R\), on the
almost-sure event \(\{R>0\}\).  Then \(X\) is distributed according to the
uniform probability measure on
\[
  S_{HS}(m,n):=\{Y\in M_{m,n}(\C):\norm Y_2=1\},
\]
and \(X\) is independent of \(R\).  Moreover
\[
  R^2=\sum_{i,\alpha}|G_{i\alpha}|^2
\]
has Gamma distribution with shape \(mn\) and scale \(1\).
\end{lemma}

\begin{proof}
Identify \(M_{m,n}(\C)\) with \(\C^{mn}\), equipped with the Euclidean norm
\(\norm{\cdot}_2\).  The density of \(G\) with respect to Lebesgue measure is
\[
  z\longmapsto \pi^{-mn}\exp(-\norm z_2^2).
\]
The origin has Gaussian measure zero, so \(R>0\) almost surely.  In polar
coordinates \(z=ru\), with \(r>0\) and \(u\in S_{HS}(m,n)\), Lebesgue measure
has the form
\[
  dz=c_{m,n}\, r^{2mn-1}\,dr\,d\sigma(u),
\]
where \(\sigma\) is the uniform probability measure on \(S_{HS}(m,n)\) and
\(c_{m,n}\) is the surface area constant.  Hence the joint density of
\((R,X)\) is proportional to
\[
  e^{-r^2}r^{2mn-1}\,dr\,d\sigma(u).
\]
This factorizes into a function of \(r\) times \(\sigma(du)\), proving both
the uniformity of \(X\) and the independence of \(X\) and \(R\).  Finally,
the change of variable \(y=r^2\) gives density proportional to
\[
  y^{mn-1}e^{-y}\,dy
\]
for \(R^2\), i.e. the Gamma distribution with shape \(mn\) and scale \(1\).
\end{proof}

\begin{lemma}[Gaussian operator norm]
\label{lem:appendixB-gaussian-op}
Let \(m,n\ge1\), and let \(H\in M_{m,n}(\C)\) have independent standard
complex Gaussian entries.  There is a universal constant \(C\) such that
\[
  \E\norm H_\infty\le C(\sqrt m+\sqrt n).
\]
\end{lemma}

\begin{proof}
For unit vectors \(u\in\C^m\) and \(v\in\C^n\),
\[
  \langle u,Hv\rangle
  =
  \sum_{i,\alpha}\overline{u_i}H_{i\alpha}v_\alpha
\]
is a standard complex Gaussian, so
\(\Prob\{|\langle u,Hv\rangle|\ge r\}=e^{-r^2}\).
Choose \(1/4\)-nets \(N_1,N_2\) in the unit spheres of \(\C^m\) and
\(\C^n\).  The elementary volume argument gives
\[
  |N_1|\le9^{2m},\qquad |N_2|\le9^{2n}.
\]
If
\[
  M=\max_{u\in N_1,\ v\in N_2}|\langle u,Hv\rangle|,
\]
then \(\norm H_\infty\le2M\).  To see this, approximate arbitrary unit
\(u,v\) by \(u_0,v_0\) in the nets and use
\[
  |\langle u,Hv\rangle|
  \le
  |\langle u_0,Hv_0\rangle|
  +\frac14\norm H_\infty+\frac14\norm H_\infty .
\]
By the union bound,
\[
  \Prob\{M\ge r\}\le
  \exp\{2(m+n)\log9-r^2\}.
\]
Integrating this tail gives
\(\E M\le C'(\sqrt m+\sqrt n)\), and hence
\(\E\norm H_\infty\le C(\sqrt m+\sqrt n)\).
\end{proof}

\begin{revlatest}
\begin{lemma}[Median bound for nonnegative random variables]
\label{lem:appendixB-nonnegative-median-bound}
Let \(Z\ge0\) be an integrable random variable and let \(M\in\R\) satisfy
\[
  \Prob\{Z\ge M\}\ge \frac12 .
\]
Then
\[
  M\le 2\E Z .
\]
In particular, every median of a nonnegative random variable is bounded above
by twice its expectation.
\end{lemma}

\begin{proof}
If \(M\le0\), the claim is immediate.  Otherwise \(M>0\), and
\[
  \E Z
  \ge
  \E\bigl(Z\mathbf 1_{\{Z\ge M\}}\bigr)
  \ge
  M\Prob\{Z\ge M\}
  \ge
  \frac M2 .
\]
\end{proof}
\end{revlatest}

\begin{proposition}[Partial-transpose Wishart expectation in the balanced window]
\label{prop:appendixB-wishart-expectation}
For every \(\lambda>0\), there are constants \(C_\lambda>0\) and
\(d_0=d_0(\lambda)\) such that the following holds.  Let \(d\ge d_0\), let
\(s\) be an integer satisfying
\[
  \frac{\lambda}{2}d^2\le s\le 2\lambda d^2 .
\]
Let \(G\in M_{d^2,s}(\C)\) have independent standard complex Gaussian
entries and put \(W=s^{-1}GG^*\).  Then
\[
  \E\norm{W^\Gamma}_\infty\le C_\lambda .
\]
\end{proposition}

\begin{proof}
Write
\[
  W^\Gamma=\diag(W^\Gamma)+Z,\qquad Z=W^\Gamma-\diag(W^\Gamma).
\]
The diagonal is unchanged by partial transpose:
\[
  (W^\Gamma)_{ii}=W_{ii}
  =
  \frac1s\sum_{\alpha=1}^s |G_{i\alpha}|^2.
\]
Thus \(\norm{\diag(W^\Gamma)}_\infty\) is the maximum of \(d^2\) averages of
\(s\) independent mean-one exponentials.  We record the Chernoff estimate
needed here.  If \(V_i=s^{-1}\sum_{\alpha=1}^s E_{i\alpha}\), with the
\(E_{i\alpha}\)'s independent exponentials of mean one, then, for every
\(u>0\),
\[
  \Prob\left\{
    V_i\ge 1+2\sqrt{u/s}+2u/s
  \right\}\le e^{-u}.
\]
Indeed, for \(\theta\in(0,1)\),
\[
  \Prob\{V_i\ge 1+\delta\}
  \le
  e^{-\theta s(1+\delta)}(1-\theta)^{-s}.
\]
Taking \(\theta=\delta/(1+\delta)\) gives
\[
  \Prob\{V_i\ge 1+\delta\}
  \le
  \exp\{-s(\delta-\log(1+\delta))\}.
\]
With \(z=\sqrt{u/s}\) and \(\delta=2z+2z^2\), the elementary inequality
\[
  \log(1+2z+2z^2)\le 2z+z^2
\]
implies \(\delta-\log(1+\delta)\ge z^2=u/s\), which proves the tail bound.

Apply this with \(u=v+2\log d\) and take a union bound over the \(d^2\)
diagonal entries.  For every \(v\ge0\),
\[
  \Prob\left\{
    \max_i V_i
    \ge
    1+2\sqrt{\frac{v+2\log d}{s}}
      +2\frac{v+2\log d}{s}
  \right\}
  \le e^{-v}.
\]
Let
\[
  a(v)=
  1+2\sqrt{\frac{v+2\log d}{s}}
    +2\frac{v+2\log d}{s}.
\]
The function \(a\) is increasing, and
\[
  a'(v)=
  \frac1{\sqrt{s}\sqrt{v+2\log d}}+\frac2s.
\]
By the layer-cake formula and the preceding tail bound,
\[
\begin{aligned}
  \E\max_{1\le i\le d^2}V_i
  &\le
  a(0)+\int_0^\infty e^{-v}a'(v)\,dv                         \\
  &\le
  1+2\sqrt{\frac{2\log d}{s}}
    +4\frac{\log d}{s}
    +\frac1{\sqrt{s}}\int_0^\infty \frac{e^{-v}}{\sqrt v}\,dv
    +\frac2s                                                    \\
  &\le
  1+C\sqrt{\frac{\log d}{s}}
    +C\frac{\log d}{s}
\end{aligned}
\]
for \(d\ge2\).
Therefore
\[
  \E\norm{\diag(W^\Gamma)}_\infty
  \le
  1+C_\lambda\frac{\sqrt{\log d}}{d}.
\]

For the off-diagonal part we use Aubrun's moment estimate,
\cite[Lemma~3.4]{Aubrun12}: for every even \(m\ge2\) there is a polynomial
\(Q\), independent of \(d\) and \(s\), such that
\[
  \E\Tr(Z^m)
  \le
  \left(\frac2s\right)^m
  (d+Q(m))^{m+2}(\sqrt s+Q(m))^m.
\]
Since \(Z\) is Hermitian and \(m\) is even,
\[
  \E\norm Z_\infty
  \le
  \bigl(\E\Tr(Z^m)\bigr)^{1/m}.
\]
Choose even \(m=m(d)\to\infty\) slowly enough that
\[
  Q(m)=o(d),\qquad \frac{\log d}{m}\to0.
\]
Taking the \(m\)-th root of Aubrun's estimate gives
\[
  \E\norm Z_\infty
  \le
  \frac2s(d+Q(m))^{1+2/m}(\sqrt s+Q(m)),
\]
which is bounded uniformly in \(d\) because
\(\lambda d^2/2\le s\le2\lambda d^2\).  The triangle inequality gives the
result.
\end{proof}

\begin{revlatest}
\begin{lemma}[Radial transfer from Wishart to the sphere]
\label{lem:appendixB-radial-transfer-pt-operator}
Let \(G\in M_{d^2,s}(\C)\) have independent standard complex Gaussian
entries, and use its polar decomposition in the real Hilbert space
\(M_{d^2,s}(\C)\) with the Hilbert--Schmidt norm.  Thus, outside the null
event \(\{G=0\}\), set
\[
  R=\norm G_2,\qquad X=\frac GR,\qquad W=\frac1s GG^*
\]
so that \(G=RX\), with \(X\in S_{HS}(d^2,s)\).  By
Lemma~\ref{lem:appendixB-gaussian-polar}, \(X\) is uniform on the
Hilbert--Schmidt sphere and independent of \(R\).  Then
\[
  \E\norm{(XX^*)^\Gamma}_\infty
  =
  \frac1{d^2}\E\norm{W^\Gamma}_\infty .
\]
\end{lemma}

\begin{proof}
By Lemma~\ref{lem:appendixB-gaussian-polar}, \(\E R^2=d^2s\).  Since
\(G=RX\), we have
\[
  W^\Gamma
  =
  \frac1s(GG^*)^\Gamma
  =
  \frac{R^2}{s}(XX^*)^\Gamma .
\]
Taking operator norms and expectations gives
\[
  \E\norm{W^\Gamma}_\infty
  =
  \frac1s\E\left(R^2\norm{(XX^*)^\Gamma}_\infty\right)
  =
  \frac{\E R^2}{s}\E\norm{(XX^*)^\Gamma}_\infty
  =
  d^2\,\E\norm{(XX^*)^\Gamma}_\infty,
\]
where the middle equality uses independence of \(R\) and \(X\).  Rearranging
proves the lemma.
\end{proof}
\end{revlatest}

\begin{proposition}[Good set in the balanced window]
\label{prop:appendixB-good-set}
For every \(\lambda>0\), there are constants \(a,b,c,C>0\) and
\(d_0=d_0(\lambda)\) such that the following holds.  Let \(d\ge d_0\), let
\(s\) be an integer satisfying \(\lambda d^2/2\le s\le2\lambda d^2\), and
let \(X\) be uniform on \(S_{HS}(d^2,s)\).  Then the set
\[
  \Omega
  =
  \left\{
    X:\norm X_\infty\le\frac a d,\quad
    \norm{(XX^*)^\Gamma}_\infty\le\frac b{d^2}
  \right\}
\]
satisfies
\[
  \Prob(\Omega^c)\le C e^{-c d^2}.
\]
\end{proposition}

\begin{proof}
Let \(G\in M_{d^2,s}(\C)\) have independent standard complex Gaussian
entries.  On the almost-sure event \(G\ne0\), put
\[
  R=\norm G_2,\qquad X=\frac GR .
\]
The polar law in Lemma~\ref{lem:appendixB-gaussian-polar} says precisely that
this angular variable \(X\) has the uniform surface law on \(S_{HS}(d^2,s)\)
and is independent of the radial variable \(R\).
First consider \(h(X)=\norm X_\infty\).  It is \(1\)-Lipschitz on the
Hilbert--Schmidt sphere.  By independence of \(R\) and \(X\),
\[
  \E\norm G_\infty=\E R\,\E\norm X_\infty.
\]
Since \(R^2\) has Gamma distribution with shape \(N=d^2s\) and scale \(1\),
\[
  \E R^2=N,\qquad \E R^4=N(N+1).
\]
The Paley--Zygmund inequality applied to \(R^2\) gives
\[
  \Prob\{R^2\ge N/2\}
  \ge
  \frac{(N/2)^2}{N(N+1)}
  \ge \frac18 .
\]
Consequently
\[
  \E R\ge c\sqrt{d^2s}.
\]
Together with Lemma~\ref{lem:appendixB-gaussian-op}, applied with
\(m=d^2\) and \(n=s\), this gives
\[
  \E\norm X_\infty\le \frac{C_\lambda}{d}.
\]
\begin{revlatest}
By Lemma~\ref{lem:appendixB-nonnegative-median-bound}, any median \(M_h\)
satisfies \(M_h\le2\E h\).  Levy's lemma in the form of
Lemma~\ref{lem:appendixB-spherical-levy}, applied on the real sphere of
dimension \(2d^2s-1\), gives, with deviation \(C_\lambda/d\),
\end{revlatest}
\[
  \Prob\left\{\norm X_\infty>\frac a d\right\}\le e^{-c d^2}.
\]

Now set \(g(X)=\norm{(XX^*)^\Gamma}_\infty\).  On the event
\(\Omega_1=\{\norm X_\infty\le a/d\}\), the map \(g\) is
\(C_\lambda/d\)-Lipschitz, because
\[
\begin{aligned}
 |g(X)-g(Y)|
 &\le
 \norm{(XX^*)^\Gamma-(YY^*)^\Gamma}_\infty        \\
 &\le
 \norm{XX^*-YY^*}_2
 \le
 (\norm X_\infty+\norm Y_\infty)\norm{X-Y}_2 .
\end{aligned}
\]
\begin{revlatest}
Furthermore, Lemma~\ref{lem:appendixB-radial-transfer-pt-operator} and
Proposition~\ref{prop:appendixB-wishart-expectation} give
\end{revlatest}
\[
  \E g(X)
  \le
  \frac{C_\lambda}{d^2}.
\]
\begin{revlatest}
Again Lemma~\ref{lem:appendixB-nonnegative-median-bound} bounds any median
\(M_g\) by \(2\E g\).  Applying the localized Levy lemma to \(g\) on
\(\Omega_1\), with deviation \(C_\lambda/d^2\), gives
\end{revlatest}
\[
  \Prob\left\{
    \norm{(XX^*)^\Gamma}_\infty>\frac b{d^2}
  \right\}
  \le
  C e^{-c d^2}.
\]
The union bound proves the proposition.
\end{proof}

\subsection{Proof of Theorem~\ref{thm:appendixB-local-concentration}}

On the good set \(\Omega\) from Proposition~\ref{prop:appendixB-good-set},
Lemma~\ref{lem:appendixB-moment-lipschitz} gives
\[
  |f_k(X)-f_k(Y)|
  \le
  \frac{C_{k,\lambda}k}{d^{2k-2}}\norm{X-Y}_2
  \qquad (X,Y\in\Omega).
\]
The real dimension of the Hilbert--Schmidt sphere is \(2d^2s-1\), which is
of order \(d^4\) and is at least \(2\) for all sufficiently large \(d\).
Applying Lemma~\ref{lem:appendixB-local-levy} with
\[
  L=\frac{C_{k,\lambda}k}{d^{2k-2}},
  \qquad
  t=\frac12\varepsilon d^{-(2k-2)},
\]
and using
\(\{|f_k-M_k|\ge \varepsilon d^{-(2k-2)}\}
  \subset
  \{|f_k-M_k|>\frac12\varepsilon d^{-(2k-2)}\}\),
we obtain, for every median \(M_k\) of \(f_k\),
\[
  \Prob\left\{
    |f_k(X)-M_k|\ge \varepsilon d^{-(2k-2)}
  \right\}
  \le
  C e^{-c d^2}
  +
  4\exp\left(-c\frac{d^4\varepsilon^2}{k^2}\right).
\]

It remains to replace the median by the expectation.  Since partial
transpose is an isometry for the Hilbert--Schmidt norm,
\[
  \norm{(XX^*)^\Gamma}_2=\norm{XX^*}_2.
\]
Moreover
\[
  \norm{XX^*}_2^2=\Tr((X^*X)^2)
  \le \Tr(X^*X)^2=\norm X_2^4=1.
\]
Thus, if \((\alpha_j)\) are the eigenvalues of \((XX^*)^\Gamma\), then
\(\sum_j|\alpha_j|^2\le1\).  Therefore \(|f_k(X)|\le1\) for \(k\ge2\),
while \(f_1(X)=1\).  From now on choose a median \(M_k\in[-1,1]\); this is
possible because \(f_k\) takes values in \([-1,1]\).  Hence
\[
  |f_k(X)-M_k|\le2
\]
for every \(X\).  This boundedness is not used to integrate a constant tail
on \([0,\infty)\).  It is used to truncate the layer-cake identity:
\[
\begin{aligned}
  |\E f_k-M_k|
  &\le \E |f_k-M_k|                                      \\
  &=\int_0^2
      \Prob\{|f_k(X)-M_k|\ge u\}\,du .
\end{aligned}
\]
For \(0<u\le2\), the localized Levy estimate applied with this value of
\(u\), together with \(n=2d^2s\asymp_\lambda d^4\) and
\(L=C_{k,\lambda}k d^{-(2k-2)}\), gives
\[
  \Prob\{|f_k(X)-M_k|\ge u\}
  \le
  C e^{-c d^2}
  +
  4\exp\left(-c\,\frac{d^{4k}u^2}{k^2}\right).
\]
Therefore
\[
\begin{aligned}
  |\E f_k-M_k|
  &\le
  2C e^{-c d^2}
  +
  4\int_0^\infty
    \exp\left(-c\,\frac{d^{4k}u^2}{k^2}\right)\,du        \\
  &\le
  C e^{-c d^2}+\frac{C_{k,\lambda}k}{d^{2k}} .
\end{aligned}
\]
If the right hand side in the last display is at most
\(\varepsilon d^{-(2k-2)}/2\), then
\[
  \{|f_k-\E f_k|\ge\varepsilon d^{-(2k-2)}\}
  \subset
  \{|f_k-M_k|\ge\tfrac12\varepsilon d^{-(2k-2)}\},
\]
and the preceding median estimate gives the result after changing constants.
If instead
\[
  |\E f_k-M_k|>\frac12\varepsilon d^{-(2k-2)},
\]
then the displayed bound on \(|\E f_k-M_k|\) forces
\(\varepsilon\le C_{k,\lambda}d^{-2}\), up to changing constants and ignoring
the exponentially small term.  In that regime
\[
  \exp\left(-c\frac{d^4\varepsilon^2}{k^2}\right)
\]
is bounded below by a positive constant depending only on \(k,\lambda\), so
the right hand side of the theorem can be made at least \(1\) by increasing
\(C\).  The probability bound is then trivial.  This proves the theorem.


% Standalone-only bibliography for compiling this wrapper outside the article.
% Do not paste this block into the main article, which already uses
% \bibliographystyle{alpha} and \bibliography{references}.
\begin{thebibliography}{2}
\bibitem{Aubrun12}
G.~Aubrun,
\emph{Partial transposition of random states and non-centered semicircular
distributions},
Random Matrices: Theory Appl. \textbf{1} (2012), 1250001.

\bibitem{AubrunSzarek2017}
G.~Aubrun and S.~J. Szarek,
\emph{Alice and Bob Meet Banach: The Interface of Asymptotic Geometric
Analysis and Quantum Information Theory},
Mathematical Surveys and Monographs, vol.~223, American Mathematical
Society, 2017.
\end{thebibliography}
\end{document}
